\documentclass{beamer}
\usetheme{Madrid}
\usecolortheme{default}

% Packages
\usepackage{amsmath, amssymb}
\usepackage{graphicx}
\usepackage{booktabs}
\usepackage{xcolor}

% Title and author
\title{Routing Strategies for Multi-Agent LLM Systems}
\subtitle{Semantic, Performance-Weighted, and LLM Meta-Routing}
\author{Zhaoheng Li}
\institute{University of Toronto \\ Department of Electrical and Computer Engineering}
\date{December 2025}

\begin{document}

% Slide 1: Title
\frame{\titlepage}

% Slide 2: Problem Statement
\begin{frame}
\frametitle{The Challenge}
\begin{itemize}
    \item Multi-agent LLM systems leverage specialized domain expertise
    \pause
    \item But they introduce a critical problem: \textbf{query routing}
    \pause
    \item Which agent should handle this query?
    \pause
    \item Trade-offs:
    \begin{itemize}
        \item \textbf{Accuracy}: Route to the right agent
        \item \textbf{Efficiency}: Minimize routing overhead
        \item \textbf{Adaptability}: Improve over time
    \end{itemize}
\end{itemize}
\end{frame}

% Slide 3: Our Four Agents
\begin{frame}
\frametitle{Specialist Agents}
\begin{columns}
\column{0.5\textwidth}
\textbf{StatisticsAgent}
\begin{itemize}
    \item Hypothesis testing
    \item Confidence intervals
    \item Experimental design
\end{itemize}

\textbf{MLAgent}
\begin{itemize}
    \item Model training
    \item Feature engineering
    \item Hyperparameter tuning
\end{itemize}

\column{0.5\textwidth}
\textbf{DataEngAgent}
\begin{itemize}
    \item SQL \& schema design
    \item ETL pipelines
    \item Data cleaning
\end{itemize}

\textbf{VizAgent}
\begin{itemize}
    \item Chart selection
    \item Dashboard design
    \item Data storytelling
\end{itemize}
\end{columns}
\end{frame}

% Slide 4: Method 1 - Semantic Similarity
\begin{frame}
\frametitle{Method 1: Semantic Similarity Baseline}
\begin{center}
    \includegraphics[width=0.6\textwidth]{method1_diagram.png}
\end{center}
\begin{itemize}
    \item Use pre-trained sentence embeddings (all-MiniLM-L6-v2)
    \item Encode query and agent descriptions
    \item Route to agent with highest cosine similarity
    \pause
    \item \textbf{Advantages}: Training-free, fast, stable
    \item \textbf{Disadvantage}: No learning from history
\end{itemize}
\end{frame}

% Slide 5: Method 2 - Performance-Weighted Routing
\begin{frame}
\frametitle{Method 2: Performance-Weighted Adaptive Routing}
\begin{equation*}
S_a(q) = 0.20 \cdot \text{semantic}(q,a) + 0.35 \cdot \text{accuracy}_a + 0.35 \cdot \text{indomain}_a + 0.10 \cdot \text{latency\_score}_a
\end{equation*}

\begin{itemize}
    \item Maintains per-agent statistics:
    \begin{itemize}
        \item Accuracy rate (from ground truth)
        \item IN\_DOMAIN rate (agent self-assessment)
        \item Average latency
    \end{itemize}
    \pause
    \item Updates via exponential moving average ($\alpha = 0.1$)
    \pause
    \item Learns during both training \textit{and} testing phases
\end{itemize}
\end{frame}

% Slide 6: Method 3 - LLM Meta-Classification
\begin{frame}
\frametitle{Method 3: LLM Meta-Classification}
\begin{itemize}
    \item Use another LLM call to route queries
    \pause
    \item Meta-classifier sees:
    \begin{itemize}
        \item All four agent descriptions
        \item The user's query
    \end{itemize}
    \pause
    \item LLM decides which agent is best
    \pause
    \item \textbf{Advantage}: Can reason about complex boundaries
    \item \textbf{Disadvantage}: Extra LLM call = high latency
\end{itemize}
\end{frame}

% Slide 7: Experimental Setup
\begin{frame}
\frametitle{Evaluation Design}
\begin{itemize}
    \item \textbf{Datasets}:
    \begin{itemize}
        \item Training: 120 labeled queries (30 per agent)
        \item Testing: 200 labeled queries
    \end{itemize}
    \pause
    \item \textbf{Reproducibility}: 5 complete trials with different random seeds (22, 32, 52, 62, 82)
    \pause
    \item \textbf{Quality Validation}: Spot-check 20 sampled responses
    \pause
    \item \textbf{Metrics}:
    \begin{itemize}
        \item Routing accuracy, in-domain rate
        \item Latency (routing + total)
        \item Answer quality (5-point scale)
    \end{itemize}
\end{itemize}
\end{frame}

% Slide 8: Results - Accuracy
\begin{frame}
\frametitle{Key Results: Routing Accuracy}
\begin{center}
\begin{tabular}{lccc}
\toprule
\textbf{Method} & \textbf{Accuracy} & \textbf{In-Domain} & \textbf{Stability} \\
\midrule
Method 1: Semantic & 69.0\% & 91.8\% & ±0.0\% \\
Method 2: Performance & 72.8\% & 94.2\% & ±1.5\% \\
Method 3: LLM & 94.7\% & 99.2\% & ±0.6\% \\
\bottomrule
\end{tabular}
\end{center}

\vspace{0.3cm}

\begin{itemize}
    \item Method 2: +3.8pp over baseline
    \item Method 3: +25.7pp over baseline (but at what cost?)
\end{itemize}
\end{frame}

% Slide 9: Results - Latency
\begin{frame}
\frametitle{Key Results: Latency \& Efficiency}
\begin{center}
\begin{tabular}{lcc}
\toprule
\textbf{Method} & \textbf{Routing (ms)} & \textbf{Total (s)} \\
\midrule
Method 1: Semantic & 11.5 & 5.49 \\
Method 2: Performance & 12.0 & 6.80 (+24\%) \\
Method 3: LLM & 2410.6 & 8.24 (+50\%) \\
\bottomrule
\end{tabular}
\end{center}

\vspace{0.3cm}

\begin{itemize}
    \item Method 3: 200× higher routing overhead!
    \pause
    \item Efficiency ratio:
    \begin{itemize}
        \item Method 2: 0.16 pp accuracy per percent latency
        \item Method 3: 0.51 pp accuracy per percent latency
    \end{itemize}
\end{itemize}
\end{frame}

% Slide 10: Results - Quality
\begin{frame}
\frametitle{Quality Validation (Spot-Check)}
\begin{center}
    \includegraphics[width=0.8\textwidth]{quality_breakdown.png}
\end{center}
\begin{itemize}
    \item Mean quality: 4.47/5.0 (90\% good or excellent)
    \item StatisticsAgent leads: 4.71/5.0
    \item Even misrouted queries maintain standards
\end{itemize}
\end{frame}

% Slide 11: The Trade-off
\begin{frame}
\frametitle{The Accuracy-Latency Trade-off}
\begin{center}
    \textbf{When to use each method?}
\end{center}

\begin{itemize}
    \item \textbf{Method 1 (69\%, 5.49s)}
    \begin{itemize}
        \item Quick prototyping, resource-constrained environments
        \item Best when latency is critical
    \end{itemize}
    \pause
    \item \textbf{Method 2 (72.8\%, 6.80s)}
    \begin{itemize}
        \item Balanced accuracy-efficiency
        \item Learns over time (improves in deployment)
        \item Best for most production systems
    \end{itemize}
    \pause
    \item \textbf{Method 3 (94.7\%, 8.24s)}
    \begin{itemize}
        \item Critical domains (medical, legal, financial)
        \item When quality trumps latency
    \end{itemize}
\end{itemize}
\end{frame}

% Slide 12: Key Insight
\begin{frame}
\frametitle{Key Innovation: Dual-Signal Routing}
\begin{itemize}
    \item \textbf{External signal}: Accuracy (ground truth)
    \begin{itemize}
        \item What the data says
    \end{itemize}
    \pause
    \item \textbf{Internal signal}: IN\_DOMAIN (agent self-assessment)
    \begin{itemize}
        \item What the agent believes about its capabilities
    \end{itemize}
    \pause
    \item \textbf{Finding}: Combining both signals (50-50 weight) works better than either alone
    \pause
    \item Method 2's 94.2\% in-domain rate shows agents are well-calibrated
    \item Enables learning without ground-truth labels during deployment!
\end{itemize}
\end{frame}

% Slide 13: Discussion
\begin{frame}
\frametitle{Why Method 2 Works}
\begin{itemize}
    \item Continuous learning: improves as it routes queries
    \pause
    \item Balanced approach: combines signals
    \begin{itemize}
        \item 20\% semantic relevance
        \item 35\% accuracy from history
        \item 35\% agent confidence
        \item 10\% latency
    \end{itemize}
    \pause
    \item Practical deployment: doesn't require perfect routing
    \begin{itemize}
        \item 72.8\% correct routing yields 4.50/5.0 quality
        \item 27.2\% misroutes still acceptable
    \end{itemize}
\end{itemize}
\end{frame}

% Slide 14: Limitations & Future Work
\begin{frame}
\frametitle{Limitations \& Future Directions}
\textbf{Current Limitations:}
\begin{itemize}
    \item Only 4 agents, 200 queries
    \item Single LLM backbone
    \item Binary in-domain labels
\end{itemize}

\pause

\textbf{Future Work:}
\begin{itemize}
    \item Hybrid routing (blend all three methods)
    \item Scale to 10-50 agent ensembles
    \item Online learning with streaming feedback
    \item Heterogeneous agents (different LLMs, APIs)
\end{itemize}
\end{frame}

% Slide 15: Conclusion
\begin{frame}
\frametitle{Conclusion}
\begin{center}
    \Large \textbf{Intelligent routing unlocks multi-agent systems}
\end{center}

\vspace{0.3cm}

\begin{itemize}
    \item Three methods span the accuracy-latency spectrum
    \pause
    \item Performance-weighted routing (Method 2) offers best balance
    \pause
    \item Dual-signal approach (accuracy + self-assessment) is key
    \pause
    \item Multi-agent systems can leverage expertise \textit{efficiently}
\end{itemize}

\vspace{0.5cm}

\begin{center}
    \Large \textcolor{blue}{\textbf{Questions?}}
\end{center}
\end{frame}

\end{document}
